\section{Metodología de Entrenamiento}

Esta sección describe la metodología completa de entrenamiento utilizada en el experimento,
incluyendo el preprocesamiento de datos, la configuración del optimizador, las técnicas de
regularización y los mecanismos de monitoreo. Todos los hiperparámetros se centralizan en
un archivo de configuración para garantizar la reproducibilidad del experimento.

% ==============================================================================
\subsection{Preprocesamiento de Datos}

El dataset utilizado es IMDB, que contiene 50,000 reseñas de películas etiquetadas como
positivas o negativas. Las 25,000 muestras originales de entrenamiento se dividen mediante
permutación aleatoria (con semilla fija \texttt{seed=42}) en 20,000 muestras de entrenamiento
y 5,000 de validación. El conjunto de prueba se mantiene intacto con 25,000 muestras
independientes del split de test de IMDB.

El pipeline de preprocesamiento de texto consiste en los siguientes pasos:
\begin{enumerate}
    \item \textbf{Remoción de etiquetas HTML:} Se eliminan todas las etiquetas HTML
          (e.g., \texttt{<br />}, \texttt{<p>}) mediante expresión regular.
    \item \textbf{Conversión a minúsculas:} Todo el texto se normaliza a minúsculas.
    \item \textbf{Eliminación de caracteres no alfabéticos:} Se conservan únicamente
          letras y espacios, removiendo puntuación y dígitos.
    \item \textbf{Tokenización por espacios:} El texto limpio se divide en tokens
          separados por espacios en blanco.
    \item \textbf{Construcción de vocabulario:} Se construye un vocabulario con un máximo
          de 25,000 palabras, considerando únicamente aquellas con frecuencia mínima de 2
          en el conjunto de entrenamiento (20K muestras). Se reservan los índices 0 y 1
          para los tokens especiales \texttt{<PAD>} y \texttt{<UNK>} respectivamente.
    \item \textbf{Padding/truncamiento:} Todas las secuencias se ajustan a una longitud
          fija de 200 tokens. Las secuencias más largas se truncan y las más cortas se
          rellenan con ceros al final.
\end{enumerate}

La función \texttt{clean\_text} implementa los pasos 1--4 del preprocesamiento:

% --- Función clean_text: preprocesamiento de texto ---
\begin{minted}{python}
def clean_text(text: str) -> str:
    """Clean a review text by removing HTML, lowercasing,
    removing punctuation, and tokenizing."""
    # Remove HTML tags
    text = re.sub(r'<[^>]+>', ' ', text)
    # Convert to lowercase
    text = text.lower()
    # Remove non-alpha characters (keep only letters and spaces)
    text = re.sub(r'[^a-z\s]', '', text)
    # Tokenize by splitting on whitespace and rejoin
    tokens = text.split()
    return ' '.join(tokens)
\end{minted}

La función \texttt{pad\_sequence} implementa el ajuste de longitud fija, truncando secuencias
largas y rellenando con ceros las secuencias cortas:

% --- Función pad_sequence: truncamiento y padding ---
\begin{minted}{python}
def pad_sequence(encoded: list[int], max_len: int) -> list[int]:
    """Pad or truncate an encoded sequence to exactly max_len.

    If the sequence is longer than max_len, truncate it.
    If shorter, pad with zeros (PAD=0) at the end.
    """
    if len(encoded) > max_len:
        return encoded[:max_len]
    elif len(encoded) < max_len:
        return encoded + [0] * (max_len - len(encoded))
    else:
        return list(encoded)
\end{minted}

% ==============================================================================
\subsection{Función de Pérdida y Optimizador}

Se utiliza \texttt{BCEWithLogitsLoss} (Binary Cross-Entropy con logits) como función de
pérdida. Esta función combina internamente una capa sigmoide con la entropía cruzada
binaria, proporcionando mayor estabilidad numérica que aplicar sigmoide y BCE por separado.
Para la clasificación final, se aplica un umbral de 0.5 sobre la salida sigmoide del logit:
una predicción se considera positiva si $\sigma(\hat{y}) > 0.5$.

El optimizador seleccionado es Adam \cite{kingma2015} con los siguientes hiperparámetros:
\begin{itemize}
    \item Learning rate: $\eta = 0.001$
    \item $\beta_1 = 0.9$ (media móvil exponencial del primer momento)
    \item $\beta_2 = 0.999$ (media móvil exponencial del segundo momento)
\end{itemize}

El siguiente fragmento muestra la configuración del criterio de pérdida y el optimizador:

% --- Configuración de criterion y optimizer ---
\begin{minted}{python}
# Loss function and optimizer
criterion = nn.BCEWithLogitsLoss()
optimizer = torch.optim.Adam(model.parameters(), lr=lr)

# Early stopping state
best_val_loss = float("inf")
epochs_without_improvement = 0

# Ensure checkpoint directory exists
os.makedirs(checkpoint_dir, exist_ok=True)

# Checkpoint path based on model's rnn_type
checkpoint_path = os.path.join(
    checkpoint_dir, f"{model.rnn_type}_best.pt"
)

# Global step counter for gradient monitoring
global_step = 0
\end{minted}

% ==============================================================================
\subsection{Bucle de Entrenamiento}

El bucle de entrenamiento implementa el flujo estándar de forward pass, cálculo de pérdida,
backward pass y actualización de parámetros. Adicionalmente, se registran las normas L2 de
los gradientes por capa durante los primeros 100 pasos globales de entrenamiento, lo que
permite analizar la estabilidad del flujo de gradientes en las etapas iniciales.

% --- Bucle de entrenamiento: forward, backward, gradient monitoring ---
\begin{minted}{python}
for inputs, labels in train_loader:
    inputs = inputs.to(device)
    labels = labels.to(device).float()

    # Forward pass
    optimizer.zero_grad()
    logits = model(inputs).squeeze(1)
    loss = criterion(logits, labels)

    # Backward pass
    loss.backward()

    # Record gradient norms for the first 100 steps
    global_step += 1
    if global_step <= GRADIENT_MONITOR_STEPS:
        for name, param in model.named_parameters():
            if param.grad is not None:
                grad_norm = param.grad.data.norm(2).item()
                if name not in history.gradient_norms:
                    history.gradient_norms[name] = []
                history.gradient_norms[name].append(grad_norm)
\end{minted}

% ==============================================================================
\subsection{Gradient Clipping}

Para prevenir la explosión de gradientes (\textit{exploding gradients}), se aplica gradient
clipping con norma L2 máxima de 5.0. Este mecanismo se ejecuta después del cálculo de
gradientes (\texttt{loss.backward()}) y antes del paso del optimizador
(\texttt{optimizer.step()}), recortando la norma global del vector de gradientes cuando
excede el umbral configurado.

% --- Gradient clipping ---
\begin{minted}{python}
    # Gradient clipping
    torch.nn.utils.clip_grad_norm_(
        model.parameters(), max_norm=clip_norm
    )

    # Optimizer step
    optimizer.step()

    # Track metrics
    running_loss += loss.item() * inputs.size(0)
    predictions = (torch.sigmoid(logits) > 0.5).long()
    correct += (predictions == labels.long()).sum().item()
    total += labels.size(0)
\end{minted}

% ==============================================================================
\subsection{Early Stopping}

Se implementa early stopping para detener el entrenamiento cuando la pérdida de validación
deja de mejorar, evitando así el sobreajuste. El mecanismo monitorea la pérdida de
validación al final de cada época: si no se observa mejora durante 5 épocas consecutivas
(\texttt{patience=5}), el entrenamiento se detiene anticipadamente. El número máximo de
épocas es 10. Cuando se detecta una mejora, se guarda un checkpoint del mejor modelo.

% --- Early stopping: comparación val_loss y guardado de checkpoint ---
\begin{minted}{python}
# --- Early stopping check ---
if epoch_val_loss < best_val_loss:
    best_val_loss = epoch_val_loss
    epochs_without_improvement = 0
    # Save best checkpoint
    torch.save(model.state_dict(), checkpoint_path)
    print(f"  Best model saved to {checkpoint_path}")
else:
    epochs_without_improvement += 1
    print(
        f"  No improvement "
        f"({epochs_without_improvement}/{patience})"
    )
    if epochs_without_improvement >= patience:
        print(
            f"  Early stopping triggered after "
            f"epoch {epoch + 1}"
        )
        break
\end{minted}

% ==============================================================================
\subsection{Reproducibilidad y Configuración}

Para garantizar la reproducibilidad completa del experimento, se fija la semilla
\texttt{seed=42} en todos los generadores de números aleatorios relevantes:
\begin{itemize}
    \item \texttt{random.seed(42)} — módulo \texttt{random} de Python
    \item \texttt{np.random.seed(42)} — NumPy
    \item \texttt{torch.manual\_seed(42)} — PyTorch CPU
    \item \texttt{torch.cuda.manual\_seed\_all(42)} — PyTorch CUDA (todas las GPUs)
    \item \texttt{torch.backends.cudnn.deterministic = True} — cuDNN determinístico
    \item \texttt{torch.backends.cudnn.benchmark = False} — desactiva autotuning
\end{itemize}

El tamaño de batch utilizado es de 64 muestras. El conjunto de prueba consta de 25,000
muestras del split de test original de IMDB, evaluado una única vez al final del
entrenamiento con el mejor checkpoint guardado. Durante los primeros 100 pasos de
entrenamiento se registran las normas L2 de los gradientes por capa, lo que permite
diagnosticar problemas de vanishing o exploding gradients en las etapas iniciales del
aprendizaje.
